\documentclass[11pt,a4paper]{article}

\usepackage[T1]{fontenc}
\usepackage[utf8]{inputenc}
\usepackage{libertinus}
\usepackage{microtype}
\usepackage[a4paper,margin=19mm,headheight=15pt]{geometry}
\usepackage{amsmath,amssymb,mathtools}
\usepackage{booktabs,longtable,tabularx,array,multirow}
\usepackage{graphicx,float,placeins,pdflscape}
\usepackage[most]{tcolorbox}
\usepackage{xcolor}
\usepackage{fancyhdr}
\usepackage{enumitem}
\usepackage[numbers,sort&compress]{natbib}
\usepackage{hyperref}

\definecolor{navy}{HTML}{163A5F}
\definecolor{teal}{HTML}{087E8B}
\definecolor{green}{HTML}{2A7F62}
\definecolor{orange}{HTML}{D66B18}
\definecolor{pale}{HTML}{EEF5F7}
\definecolor{warm}{HTML}{FFF6E8}
\definecolor{redpale}{HTML}{FFF0F0}
\definecolor{muted}{HTML}{4B5563}
\hypersetup{
  colorlinks=true,
  linkcolor=navy,
  citecolor=teal,
  urlcolor=teal,
  pdfauthor={DMD initial validation implementation},
  pdftitle={Initial validation results for DMD on sparse calcium-event data}
}
\pagestyle{fancy}
\fancyhf{}
\lhead{DMD initial validation results}
\rhead{Accepted run --- 16 July 2026}
\cfoot{\thepage}
\setlength{\parindent}{0pt}
\setlength{\parskip}{0.48em}
\setlist{nosep,leftmargin=1.6em}
\renewcommand{\arraystretch}{1.14}
\newcolumntype{Y}{>{\raggedright\arraybackslash}X}
\newcolumntype{P}[1]{>{\raggedright\arraybackslash}p{#1}}
\newcolumntype{C}[1]{>{\centering\arraybackslash}p{#1}}

\newtcolorbox{resultbox}[1][]{breakable,colback=pale,colframe=teal,
  boxrule=0.8pt,arc=2pt,left=7pt,right=7pt,top=6pt,bottom=6pt,
  title=#1,fonttitle=\bfseries}
\newtcolorbox{passbox}[1][]{breakable,colback=green!5,colframe=green,
  boxrule=0.8pt,arc=2pt,left=7pt,right=7pt,top=6pt,bottom=6pt,
  title=#1,fonttitle=\bfseries}
\newtcolorbox{failbox}[1][]{breakable,colback=redpale,colframe=orange,
  boxrule=0.8pt,arc=2pt,left=7pt,right=7pt,top=6pt,bottom=6pt,
  title=#1,fonttitle=\bfseries}
\newtcolorbox{cautionbox}[1][]{breakable,colback=warm,colframe=orange,
  boxrule=0.8pt,arc=2pt,left=7pt,right=7pt,top=6pt,bottom=6pt,
  title=#1,fonttitle=\bfseries}

\newcommand{\T}{^{\mathsf T}}
\newcommand{\R}{\mathbb{R}}
\DeclareMathOperator{\tr}{tr}
\newcommand{\runid}{20260716T233537+0900\_ca0ddd9\_3dc26e9ab8}
\newcommand{\runroot}{../results/03_initial_validation/runs/20260716T233537+0900_ca0ddd9_3dc26e9ab8}
\newcommand{\figfull}[2]{%
  \begin{figure}[p]
    \centering
    \includegraphics[width=\textwidth,height=0.78\textheight,keepaspectratio]{#1}
    \caption{#2}
  \end{figure}}

\begin{document}
\hypersetup{pageanchor=false}

\begin{titlepage}
  \centering
  \vspace*{1.2cm}
  {\color{navy}\rule{\textwidth}{1.2pt}}\\[0.9cm]
  {\Huge\bfseries Initial validation of DMD\\[0.15em]
  on sparse calcium-event data\par}
  \vspace{0.45cm}
  {\Large Step-by-step implementation, diagnostics, controls, and gate decision\par}
  \vspace{0.8cm}
  {\color{navy}\rule{\textwidth}{0.6pt}}\\[0.8cm]

  \begin{tcolorbox}[colback=redpale,colframe=orange,width=0.94\textwidth,arc=2pt]
  \centering
  {\Large\bfseries Frozen screening decision: FAIL}\\[0.4em]
  The native deconvolved arm predicts better than persistence, but the pilot
  fails known-system recovery, invariant-subspace reproducibility, tracking
  resolution, and the coordinated-mode null qualifier. Intensive sliding-window
  DMD/Grassmann tracking is not authorized from this checkpoint.
  \end{tcolorbox}

  \vfill
  \begin{tabular}{rl}
    Project: & \texttt{ca\_imaging\_network\_analysis/DMD}\\
    Report: & Version 1.0\\
    Accepted run: & \texttt{\runid}\\
    Date: & 16 July 2026\\
    Protocol: & \texttt{03\_dmd\_initial\_validation\_plan.pdf}
  \end{tabular}
  \vfill
  {\small Private DMD research artifact; not part of the summer-school tutorial.}
\end{titlepage}

\hypersetup{pageanchor=true}
\pagenumbering{roman}
\tableofcontents
\clearpage
\pagenumbering{arabic}

\section{Executive conclusion}

\begin{failbox}[Decision in one sentence]
The intended 300-frame sliding-window DMD analysis is \textbf{not presently a
defensible basis for tracking brain-state dynamics}: prediction exists, but the
dynamic modes are neither recoverable under matched sparse observations nor
reproducible/resolved well enough for Grassmann-manifold tracking.
\end{failbox}

This conclusion is deliberately narrower than ``DMD does not work.'' The
implementation is numerically healthy, reproduces the direct precedent's
operator, and recovers continuously observed known systems. The failure occurs
after the latent dynamics are observed as extremely sparse deconvolved events or
calcium-like traces. In the empirical data, a heavily regularized DMD model often
beats persistence, but independent per-PC autoregression explains almost all of
that advantage. Neuronwise negative controls retain similar predictability.

The five frozen gates gave:

\begin{center}
\begin{tabular}{@{}P{0.37\textwidth}C{0.14\textwidth}P{0.38\textwidth}@{}}
\toprule
Gate & Result & Decisive observation\\
\midrule
1. Geometry and numerical health & \textcolor{green}{\textbf{PASS}} &
No illegal joins; 100\% finite; no explosive forecasts.\\
2. Known-system recovery & \textcolor{orange}{\textbf{FAIL}} &
Classification 0.55 and median subspace overlap 0.430.\\
3. Empirical prediction & \textcolor{green}{\textbf{PASS}} &
Conditional P-arm lower bootstrap bound versus persistence above zero in 4/4
fixed cells.\\
4. Empirical reproducibility & \textcolor{orange}{\textbf{FAIL}} &
0/4 cells reach $S_{\rm sub}\geq0.80$; matching failures occur.\\
5. Tracking resolution & \textcolor{orange}{\textbf{FAIL}} &
$R_{\rm track}=0.819$ and $0.854$, versus required $<0.5$.\\
\bottomrule
\end{tabular}
\end{center}

No cell satisfied the additional requirement that full-DMD gain over diagonal
AR be both positive and above the 95th percentile of the neuronwise
circular-shift null. Therefore the report supports \emph{forecastable marginal
relaxation} at most, not coordinated or rotational population modes.

\begin{resultbox}[Recommended pause]
Do not proceed automatically to all-recording sliding-window DMD or Germain-style
Grassmann tracking. The user review should decide whether to stop this route or
commission a new plan for a calcium-aware latent state-space model, coarser
rate/count observations, or much longer windows. Any revision is a new protocol,
not a reinterpretation of this failed screen.
\end{resultbox}

\figfull{\runroot/07_decision/20_gate_summary.png}{Final frozen-gate summary. Green marks the two thresholds that passed; orange marks the three that failed.}
\FloatBarrier

\section{Scope and interpretation rules}

The pilot tested whether DMD could be a credible measurement tool before asking
whether awake, NREM, or anaesthesia differ. It did \emph{not} estimate state
effects. State labels were used only to construct legal constant-label windows
and to summarize untouched results after development-only model selection.

The direct methodological anchor was the Pachitariu et al. analysis of
neuronwise RMS-standardized deconvolved activity with lagged ridge DMD
\citep{pachitariu2026critical}. Standard DMD background follows
\citet{schmid2010} and \citet{tu2014}. The present pilot adds frozen
recording-level PCA, contiguous held-out prediction, diagonal-AR baselines,
known-system observation controls, block bootstrap, proper-subspace projectors,
equal-neuron subsets, and neuronwise temporal nulls.

Three distinctions govern every interpretation:

\begin{enumerate}
  \item Beating persistence is evidence of predictability, not necessarily of a
  multivariate population operator.
  \item Beating diagonal AR is required for coordinated cross-PC dynamics.
  \item A smooth-looking subspace trajectory is not trackable unless
  within-window estimation noise is materially smaller than ordinary
  between-window variation.
\end{enumerate}

\section{Audit trail and implementation revisions}

\subsection{Chronology}

Intermediate results were inspected before the next expensive stage. Earlier
runs were preserved rather than overwritten:

\begin{longtable}{@{}P{0.31\textwidth}P{0.18\textwidth}P{0.42\textwidth}@{}}
\toprule
Run & Last completed stage & Purpose / reason superseded\\
\midrule
\endfirsthead
\toprule
Run & Last completed stage & Purpose / reason superseded\\
\midrule
\endhead
\texttt{20260716T224117\ldots} & Data & First figure geometry check; title/legend layout revised.\\
\texttt{20260716T224223\ldots} & Preprocessing & PC-trace layout audit; figure spacing revised.\\
\texttt{20260716T224405\ldots} & Forecast & First full predictive audit; exposed split, long-baseline, and provenance issues.\\
\texttt{20260716T232145\ldots} & Nulls & Corrected implementation audit; exposed the need to require positive gain in the null qualifier.\\
\path{20260716T233537+0900_ca0ddd9_3dc26e9ab8} & Decision & Final accepted,
source-consistent run reported here.\\
\bottomrule
\end{longtable}

\subsection{Declared amendments}

Four clarifications were frozen before the first empirical DMD fit: (i) require
$q\geq8$ for a $k=4$ tracking subspace, (ii) use exact global maximin window
placement with earliest ties, (iii) make development validation local 80/20 to
match deployment, and (iv) add a seven-frame stored-smoothing edge sensitivity.

The stepwise audit then produced three explicit corrections before the accepted
rerun:

\begin{enumerate}
  \item Gate 2 failed, so downstream stages continued only under the configured
  label \path{diagnostic_completion_only_not_confirmatory_validation}.
  \item The edge-trimmed sensitivity now preserves the original frame-240
  train/test boundary: after retaining frames $[7,293)$, the trimmed-coordinate
  fitting stop is 233. The earlier implementation had inadvertently moved it.
  \item The long-block diagonal AR ridge is selected independently within the
  nested development split. Fourteen of 20 arm-by-cell fits choose a ridge
  different from the full DMD model.
\end{enumerate}

The final null qualifier also requires a \emph{positive} full-over-diagonal gain,
not merely a value above a negative null percentile. Incomplete eigengroup
matching automatically makes a cell ineligible for Gates 4 and 5.

\section{Data, windows, and leakage safeguards}

\subsection{Recording inventory}

\begin{center}
\begin{tabular}{@{}lrrrrl@{}}
\toprule
Recording & Neurons & Frames & Minutes & Segments & Role\\
\midrule
\texttt{example\_data} & 1,000 & 9,350 & 20.37 & 1 & precedent check\\
\texttt{mouse02\_sleep} & 6,574 & 18,700 & 40.74 & 2 & awake / NREM\\
\texttt{mouse07\_ane} & 4,337 & 19,000 & 41.39 & 2 & awake / anaesthesia\\
\bottomrule
\end{tabular}
\end{center}

At 7.65 Hz, a 300-frame window lasts 39.22 s. Native-arm models fit the first
240 frames (31.37 s) and score the final 60 frames (7.84 s). The prespecified
1,500-frame diagnostic lasts 196.08 s and uses a nested 64/16/20 split.

Six 300-frame windows were selected in each recording--label cell by globally
maximizing minimum start separation. Chronological positions 1, 3, and 5 are
development; 2, 4, and 6 are untouched evaluation. No published
``used-frame'' mask was applied, so the early anaesthesia interval remains
available. This is appropriate for a general dataset-use tutorial and avoids
silently deleting transition-related frames.

\figfull{\runroot/00_data/00_window_geometry.png}{Raw state labels, exact development/evaluation windows, acquisition boundaries, and prespecified long blocks before any DMD fit.}

\begin{passbox}[Geometry audit]
There are 24 deployment windows, 12 development windows, 12 untouched
evaluation windows, and four long diagnostics. There are zero development--
evaluation overlaps, zero state-label violations, and zero acquisition-boundary
crossings. Snapshot pairs are formed separately inside each legal window or
bootstrap block. Thirteen deployment windows touch the approximate support of
the stored 15-frame smoother; this motivates the explicit S$z$/S$c$ sensitivity,
but it does not affect the unsmoothed P arm.
\end{passbox}

\section{Stage 1: candidate observations and frozen PCA}

Five observation arms were fit from development data only:

\begin{center}
\begin{tabular}{@{}llP{0.48\textwidth}@{}}
\toprule
Arm & Signal / scaling & Purpose\\
\midrule
P & deconvolved; RMS+$10^{-3}$ & Pachitariu-aligned anchor\\
B4 & four-frame event mass; RMS+$10^{-3}$ & coarser mean-rate/count-like arm\\
S$z$ & stored smoothed; RMS+$10^{-3}$ & code-matched smoothing control\\
S$c$ & stored smoothed; center only & scaling sensitivity\\
F & $\Delta F/F$; RMS+$10^{-3}$ & continuous calcium control\\
\bottomrule
\end{tabular}
\end{center}

\begin{center}
\begin{tabular}{@{}llrrr@{}}
\toprule
Recording & Arm & Exact zeros & Eligible neurons & Lag-1 autocorrelation\\
\midrule
\multirow{5}{*}{mouse02 sleep}
& P  & 98.91\% & 6,553 & 0.125\\
& B4 & 96.69\% & 6,553 & 0.249\\
& S$z$ & 90.38\% & 6,556 & 0.979\\
& S$c$ & 90.38\% & 6,556 & 0.981\\
& F  & 0\% & 6,574 & 0.699\\
\midrule
\multirow{5}{*}{mouse07 anaesthesia}
& P  & 99.51\% & 4,193 & 0.046\\
& B4 & 98.31\% & 4,193 & 0.174\\
& S$z$ & 94.94\% & 4,200 & 0.978\\
& S$c$ & 94.94\% & 4,200 & 0.979\\
& F  & 0\% & 4,337 & 0.697\\
\bottomrule
\end{tabular}
\end{center}

The event observations are therefore more discontinuous than a simple binary
description suggests: they are almost always exactly zero, with sparse
continuous-amplitude events. Stored smoothing and $\Delta F/F$ create strong
memory, so good persistence forecasts from those arms are not by themselves
evidence for dynamic modes.

\figfull{\runroot/01_preprocessing/01_signal_diagnostics.png}{Sparsity and pooled within-window autocorrelation for all five candidate observation arms.}
\figfull{\runroot/01_preprocessing/02_pca_diagnostics.png}{Cumulative development variance in the frozen recording-level PCA bases. The same basis is used for every state and evaluation window of a recording/arm.}

All ten recording--arm PCA bases passed the orthogonality check; the maximum
Frobenius error was $9.88\times10^{-15}$. No nonfinite neuron was found. The
heatmaps below include \emph{all} eligible neurons; no 1,000-neuron display cap
or row sampling is used.

\begin{landscape}
\begin{figure}[p]
\centering
\includegraphics[width=0.98\linewidth,height=0.82\textheight,keepaspectratio]{\runroot/01_preprocessing/03_all_neuron_heatmaps__mouse02_sleep.png}
\caption{Mouse02 representative untouched windows: every eligible neuron, activity ordered separately within each panel.}
\end{figure}
\end{landscape}

\begin{landscape}
\begin{figure}[p]
\centering
\includegraphics[width=0.98\linewidth,height=0.82\textheight,keepaspectratio]{\runroot/01_preprocessing/03_all_neuron_heatmaps__mouse07_ane.png}
\caption{Mouse07 representative untouched windows: every eligible neuron, including awake and anaesthesia.}
\end{figure}
\end{landscape}

\section{Stage 2A: direct-precedent implementation check}

The precedent smoke test retains all 1,000 standardized neuron dimensions in
\texttt{example\_data}, uses a two-frame lag and raw ridge $\alpha=0.01$, and
fits 9,348 legal snapshot pairs. The stable thin-SVD operator matches the
code-style Gram solve with relative Frobenius error
$1.52\times10^{-14}$. Its spectral radius is 0.658. The implementation is thus
algebraically faithful to the anchor calculation.

This is not a biological validation: the in-sample normalized residual is
0.897, and the eigenspectrum alone does not demonstrate held-out utility. The
audit also retains the precedent discrepancy between the released
$|\lambda|>0.25$ code filter (577 modes) and a real-part threshold described in
the caption (421 modes).

\figfull{\runroot/02_precedent/05_precedent_eigenplane.png}{Near-literal precedent check. Rotation summaries count the positive-imaginary member of each conjugate pair once.}

\section{Stage 2B: known-system recovery}

Two six-dimensional stable systems---real relaxation and one containing a known
rotational pair---were observed through four pipelines over 20 deterministic
seeds each. Continuous observations provide a software upper bound; event,
smoothed-event, and calcium observations test identifiability under the present
measurement problem.

\begin{center}
\scriptsize
\begin{tabular}{@{}llrrrr@{}}
\toprule
System & Observation & Class accuracy & Subspace overlap & Eigenvalue error & Skill vs persistence\\
\midrule
Real & continuous & 1.00 & 0.963 & 0.027 & 0.185\\
Real & event & 0.40 & 0.362 & 0.972 & 0.456\\
Real & smoothed event & 1.00 & 0.443 & 0.079 & 0.130\\
Real & calcium & 1.00 & 0.512 & 0.102 & 0.313\\
Rotational & continuous & 1.00 & 0.998 & 0.017 & 0.656\\
Rotational & event & 0.90 & 0.367 & 0.949 & 0.484\\
Rotational & smoothed event & 0.00 & 0.391 & 0.077 & 0.143\\
Rotational & calcium & 0.00 & 0.674 & 0.227 & 0.351\\
\bottomrule
\end{tabular}
\end{center}

Across the transformed observations, classification accuracy is 0.55 and median
subspace overlap is 0.430, below the frozen 0.80 and 0.75 thresholds.
Eigenvalue error (0.160) and forecast skill (0.324) pass their components. The
combination is scientifically important: smoothing/calcium can make a signal
easy to forecast while erasing the distinction between real and rotational
dynamics. Forecasting cannot substitute for mode recovery.

The aggregate eigenvalue error also conceals an observation-specific failure:
event observations have median error 0.954, whereas the smoother observation
classes pull the transformed-observation aggregate below threshold. The valid
conclusion is therefore that this observe--PCA--ridge-DMD pipeline does not
preserve identifiable modes through the tested sparse/calcium-like observation
maps. It is not evidence that DMD is numerically broken or that no event-data
DMD method could work.

\figfull{\runroot/03_simulation/06_known_system_recovery.png}{Known-system recovery at four observation levels. Dashed lines show frozen thresholds. Continuous observations validate the software; transformed observations expose the identifiability failure.}

\begin{failbox}[Gate 2 failure and diagnostic override]
Gate 2 failed before empirical forecasting. Downstream work was nevertheless
completed because the user requested a full smoke-test diagnosis. Every later
result is labelled diagnostic and cannot reverse this frozen failure.
\end{failbox}

\section{Stage 3A: anonymous development selection}

For rank $q$, lag $\ell$, PC score $z_t$, and relative ridge $\eta$, the fitted
operator is
\begin{equation}
  z_{t+\ell}=Bz_t+\varepsilon_t,
  \qquad
  \widehat B=YX\T(XX\T+\alpha I)^{-1},
  \qquad
  \alpha=\eta\frac{\tr(XX\T)}{q}.
\end{equation}
It is computed with a thin SVD rather than an explicit inverse. Development
selection fits the first 80\% and scores the last 20\% of each of six anonymous
development windows per recording. Labels do not enter the ranking.

Median development scores selected P globally (0.389), followed by B4 (0.293),
F (0.214), S$c$ (0.151), and S$z$ (0.128). The chosen P configurations are:

\begin{center}
\begin{tabular}{@{}llrrrr@{}}
\toprule
Recording & Role & $q$ & $\ell$ & $\eta$ & diagonal $\eta$\\
\midrule
mouse02 & prediction & 4 & 4 & 1 & $10^{-2}$\\
mouse02 & tracking & 8 & 2 & $10^{-1}$ & $10^{-4}$\\
mouse07 & prediction/tracking & 16 & 8 & 1 & 1\\
\bottomrule
\end{tabular}
\end{center}

The mouse02 rank-4 winner is prediction-only: a $k=4$ projector in $q=4$ would
be the identity. Tracking therefore uses the prospectively selected $q=8$
configuration. Several winners lie at grid boundaries, so the selected ranks
and lags should not be interpreted as precisely localized biological scales.

\figfull{\runroot/04_forecast/07_development_selection.png}{Anonymous blocked development selection. Hollow circles identify prediction winners; the vertical line marks the minimum rank eligible for a proper $k=4$ tracking subspace.}

\section{Stage 3B--G: untouched empirical forecasts and spectra}

\subsection{Prediction against persistence and diagonal AR}

The P arm's median untouched near-1 s skills are:

\begin{center}
\begin{tabular}{@{}lrrr@{}}
\toprule
Cell & vs persistence & vs diagonal AR & equal-PC vs persistence\\
\midrule
mouse02 awake & 0.449 & $-0.008$ & 0.494\\
mouse02 NREM & 0.185 & $-0.007$ & 0.044\\
mouse07 anaesthesia & 0.455 & 0.082 & 0.440\\
mouse07 awake & 0.516 & 0.009 & 0.511\\
\bottomrule
\end{tabular}
\end{center}

At near 2 s, P remains positive versus persistence in all cells (0.224--0.475),
but full-over-diagonal gains remain $-0.068$, $-0.001$, 0.004, and 0.089.
All primary operators, eigenvalues, SSE values, and forecasts are finite; the
explosive fraction is zero and the maximum spectral radius is below one.

\figfull{\runroot/04_forecast/08_heldout_forecast_skill.png}{Held-out skill over persistence at all representable physical horizons. Bars are medians and dots retain the three untouched window values.}

\begin{landscape}
\begin{figure}[p]
\centering
\includegraphics[width=0.98\linewidth,height=0.82\textheight,keepaspectratio]{\runroot/04_forecast/09_representative_predictions.png}
\caption{Chronologically earliest untouched-window P forecasts near 1 s. Strong shrinkage toward a near-flat trajectory can beat sparse-signal persistence without reproducing event peaks.}
\end{figure}
\end{landscape}

\figfull{\runroot/04_forecast/12_multivariate_gain.png}{Full-DMD skill over the independently tuned diagonal-AR baseline. The P-arm bars are near zero in three cells and modestly positive only in anaesthesia.}

\begin{cautionbox}[What the predictive pass means]
P meets the frozen conditional persistence criterion, but this is mostly not a
multivariate DMD effect.
Persistence carries a sparse event-like PC state forward; a regularized model
that shrinks toward zero can outperform it. The diagonal AR result and the
representative trajectories forbid interpreting Gate 3 as evidence for
coordinated modes. Relative to the training-mean predictor, the cell-median
near-1 s skill is positive only for mouse02 awake (0.130) and is negative for
mouse02 NREM ($-0.040$), mouse07 anaesthesia ($-0.100$), and mouse07 awake
($-0.086$). Thus Gate 3 is a conditional estimator-perturbation pass against
persistence in these fixed windows, not general predictive validity.
\end{cautionbox}

\subsection{Discrete spectra}

All selected predictive spectra are plotted directly in the discrete plane.
The P arm has zero interpretable held-out rotational pairs after requiring
stability, modulus $>0.25$, and at least three cycles in the fitting window.
Complex spectra in other arms, especially $\Delta F/F$, are not sufficient to
override their development ranking, baseline performance, or the known-system
failure.

\begin{landscape}
\begin{figure}[p]
\centering
\includegraphics[width=0.98\linewidth,height=0.82\textheight,keepaspectratio]{\runroot/04_forecast/10_heldout_eigenplanes.png}
\caption{Selected predictive-configuration spectra in all untouched windows. The unit circle is shown; points are discrete-lag eigenvalues.}
\end{figure}
\end{landscape}

\subsection{Long-block diagnostic}

The prespecified long blocks improve P's persistence skill most strongly in
mouse02 NREM, from 0.185 to 0.443. However, these fits reselect larger
cell-specific models (often $q=32$) and therefore combine more temporal support
with a larger search space. Their full-over-diagonal gains remain only 0.009 to
0.066 after independent diagonal-ridge tuning.

\begin{center}
\begin{tabular}{@{}lrrrrr@{}}
\toprule
Cell & W=300 skill & W=1500 skill & Long $q$ & Long $\ell$ & Long vs diagonal\\
\midrule
mouse02 awake & 0.449 & 0.482 & 32 & 8 & 0.012\\
mouse02 NREM & 0.185 & 0.443 & 32 & 4 & 0.023\\
mouse07 anaesthesia & 0.455 & 0.471 & 32 & 4 & 0.066\\
mouse07 awake & 0.516 & 0.450 & 16 & 8 & 0.009\\
\bottomrule
\end{tabular}
\end{center}

\figfull{\runroot/04_forecast/11_long_block_diagnostic.png}{Cell-level comparison of 300-frame and separately selected 1,500-frame models. Points above the diagonal improve with the longer diagnostic; this is not a pure window-length contrast.}

\subsection{Stored-smoothing boundary sensitivity}

After trimming seven frames at both edges while preserving the original split,
median near-1 s persistence skill changes as follows:

\begin{center}
\begin{tabular}{@{}lrrrr@{}}
\toprule
Cell & S$z$ primary & S$z$ trimmed & S$c$ primary & S$c$ trimmed\\
\midrule
mouse02 awake & 0.168 & 0.160 & 0.335 & 0.336\\
mouse02 NREM & $-0.058$ & 0.076 & $-0.269$ & 0.121\\
mouse07 anaesthesia & 0.271 & 0.263 & 0.150 & 0.177\\
mouse07 awake & 0.268 & 0.308 & 0.370 & 0.340\\
\bottomrule
\end{tabular}
\end{center}

The sign reversals in mouse02 NREM make stored-smoothed conclusions
boundary-sensitive. The full-range plot retains severe individual-window
failures rather than clipping them; the table supplies the robust cell medians.

\figfull{\runroot/04_forecast/13_smoothing_trim_sensitivity.png}{Primary versus fixed-split edge-trimmed skill for every untouched S$z$/S$c$ window. Extreme negative values are deliberately retained.}

\section{Stage 4: bootstrap prediction and mode reproducibility}

\subsection{Moving-block forecast uncertainty}

Each evaluation-window fitting segment was resampled 100 times in moving blocks
of at least 15 native frames. Pairs were created only inside blocks. Repetition
$b$ was synchronized across the three evaluation windows, their cell median was
taken, and percentile intervals were then computed. These intervals quantify
operator-estimation uncertainty for the fixed selected windows; they are not
population-level confidence intervals over animals. They condition on the
frozen PCA basis, hyperparameters, windows, and test tails and do not include
test-origin, window-selection, neuron-sampling, preprocessing-choice, or
model-selection uncertainty.

\begin{center}
\scriptsize
\begin{tabular}{@{}lrrr@{}}
\toprule
Cell & 1-s bootstrap median & 95\% interval & full-over-diagonal 95\% interval\\
\midrule
mouse02 awake & 0.450 & [0.445, 0.452] & [$-0.021$, 0.001]\\
mouse02 NREM & 0.183 & [0.169, 0.196] & [$-0.033$, 0.026]\\
mouse07 anaesthesia & 0.452 & [0.443, 0.460] & [0.035, 0.091]\\
mouse07 awake & 0.510 & [0.497, 0.519] & [$-0.030$, 0.011]\\
\bottomrule
\end{tabular}
\end{center}

Thus the conditional interval excludes zero versus persistence for all four
fixed cells, while only anaesthesia has a positive full-over-diagonal interval.

\figfull{\runroot/05_stability/14_bootstrap_forecast_intervals.png}{Synchronized moving-block intervals for P-arm skill over persistence at the near-1 s and near-2 s gates.}

\subsection{Fixed-reference invariant subspaces}

For each recording, a development-only pooled operator defines $k=4$ real
invariant dimensions ranked by mean modal-coefficient energy. Conjugate pairs
are never split. Evaluation and bootstrap eigengroups are matched to these fixed
groups; no evaluation fit is allowed to re-rank a convenient subspace.

This frozen target itself requires caution. Modal-coefficient energy is based on
the pseudoinverse of the eigenvector matrix and can be ill-conditioned or
non-additive for nonnormal/repeated modes. Mouse07's selected target comprises
two complex pairs with $|\lambda|\simeq0.077$ and 0.196, both below the
protocol's $|\lambda|>0.25$ interpretability cutoff and rapidly attenuated at
the selected eight-frame lag. Mouse02 also includes modes at 0.242 and 0.283.
Low stability therefore characterizes this prespecified tracking object; it is
not proof that no other dynamical representation exists.

Projector similarity and distance are
\begin{equation}
 S_{\rm sub}=\frac{1}{k}\tr(PP^{(b)}),
 \qquad
 d_P=\sqrt{\max(0,1-S_{\rm sub})}.
\end{equation}

\begin{center}
\scriptsize
\begin{tabular}{@{}lrrrrl@{}}
\toprule
Cell & Median $S_{\rm sub}$ & Within $d_P$ & Match fraction & Min. windows/rep & Gate eligible\\
\midrule
mouse02 awake & 0.535 & 0.549 & 0.670 & 1 & no\\
mouse02 NREM & 0.500 & 0.570 & 0.813 & 1 & no\\
mouse07 anaesthesia & 0.238 & 0.754 & 1.000 & 3 & yes\\
mouse07 awake & 0.262 & 0.753 & 1.000 & 3 & yes\\
\bottomrule
\end{tabular}
\end{center}

Mouse02's development reference selects four real one-dimensional groups, but
some evaluation/bootstrap fits resolve complex pairs instead. The
dimension-compatible matcher correctly refuses to coerce them. There are 157
recorded original/bootstrap matching failures. The two fully eligible mouse07
cells remain far below 0.80, so no cell passes the reproducibility threshold.
Reported summaries condition on successful matches and are mildly optimistic:
the assignment cost itself includes a 0.25-weighted version of the reported
subspace-overlap endpoint.

\figfull{\runroot/05_stability/15_subspace_stability.png}{Similarity to the fixed development reference for original evaluation fits and moving-block fits. Unmatched fits are recorded as failures and excluded from boxes, making the table's eligibility columns essential.}

\subsection{Tracking resolution}

The future tracking criterion compares within-window estimation distance with
ordinary same-label between-window distance:
\begin{equation}
 R_{\rm track}=\frac{\operatorname{median}(d_P^{\rm bootstrap})}
 {\operatorname{median}(d_P^{\rm between})}.
\end{equation}

Mouse02 has $R_{\rm track}=0.819$ but is ineligible because only four of six
original evaluation projectors match. Mouse07 is eligible and has
$R_{\rm track}=0.854$. Both exceed the 0.5 ceiling. Estimation noise is nearly as
large as ordinary window-to-window change; a visually smooth Grassmann path
would not be trustworthy. The between-window distances (approximately 0.685
and 0.879) are also close to independent-random-subspace expectations for the
respective $(q,k)$ dimensions (approximately 0.707 and 0.866).

\figfull{\runroot/05_stability/16_tracking_resolution.png}{Tracking-resolution components and ratios. The dashed line is the frozen $R_{\rm track}<0.5$ threshold.}

\subsection{Two disjoint neuron subsets}

Two deterministic, disjoint 2,000-neuron subsets were independently reduced to
the selected tracking rank. Their majority qualitative classes agree in all
cells: resolved rotation for mouse02 awake and no resolved rotation for the
other three cells. Agreement is useful but insufficient: it does not fix low
reference similarity or poor tracking resolution.

\figfull{\runroot/05_stability/17_neuron_subset_classes.png}{Qualitative spectral class in two equal-size disjoint neuron subsets.}

\section{Stage 5: empirical negative controls}

For every evaluation window, 100 neuronwise circular shifts were generated
inside the same acquisition-split, constant-label bout, with at least a
15-frame phase displacement. They preserve each neuron's empirical marginal and
autocorrelation approximately while disrupting synchrony. A second set of 100
stationary controls draws independently from each neuron's within-bout empirical
marginal. Both pass through the same frozen scaling, PCA, selected predictive
model, tracking model, and scoring pipeline.

Circular shifts remain predictively useful versus persistence: their median
skills are 0.320, 0.366, 0.456, and 0.457 across the four cells. Stationary-IID
controls are similarly positive (0.331, 0.370, 0.467, and 0.468). This directly
demonstrates that the empirical persistence result does not require preserved
population synchrony or temporal dynamics; on a 99.2\%-zero signal, shrinkage
toward equilibrium can by itself beat persistence.

\begin{center}
\scriptsize
\begin{tabular}{@{}lrrrrl@{}}
\toprule
Cell & Observed full/diag & Circular median & Circular 95th & Positive? & Qualifier\\
\midrule
mouse02 awake & $-0.008$ & $-0.089$ & $-0.028$ & no & fail\\
mouse02 NREM & $-0.007$ & $-0.037$ & $-0.006$ & no & fail\\
mouse07 anaesthesia & 0.082 & 0.188 & 0.235 & yes & fail\\
mouse07 awake & 0.009 & 0.023 & 0.040 & yes & fail\\
\bottomrule
\end{tabular}
\end{center}

Mouse02 awake is above a negative null percentile, but its observed gain is
itself negative and therefore cannot qualify. In the two positive cells, the
observed value is below the circular 95th percentile. Zero of four cells pass.

\begin{landscape}
\begin{figure}[p]
\centering
\includegraphics[width=0.98\linewidth,height=0.82\textheight,keepaspectratio]{\runroot/06_nulls/18_null_multivariate_gain.png}
\caption{Observed full-over-diagonal skill (orange) against synchronized cell-level circular-shift and stationary-null distributions.}
\end{figure}
\end{landscape}

Null subspace similarity is also comparable to the already-low observed value,
especially in mouse02. Neither observed nor null distributions approach the
Gate-4 threshold.

\figfull{\runroot/06_nulls/19_null_subspace_similarity.png}{Similarity to the fixed development reference under circular and stationary controls. Orange lines are original evaluation medians; dashed lines mark 0.80.}

\section{Final gate application}

\begin{longtable}{@{}P{0.18\textwidth}P{0.25\textwidth}P{0.34\textwidth}C{0.11\textwidth}@{}}
\toprule
Gate & Frozen threshold & Observed & Result\\
\midrule
\endfirsthead
\toprule
Gate & Frozen threshold & Observed & Result\\
\midrule
\endhead
1. Geometry / numerics & Zero illegal joins; at least 95\% finite &
Zero illegal geometry; 100\% finite; zero bootstrap joins; zero explosions &
\textcolor{green}{PASS}\\
2. Known systems & Class $\geq0.80$; $S_{\rm sub}\geq0.75$; eigenvalue error $\leq0.25$; skill $>0$ &
0.55; 0.430; 0.160; 0.324 & \textcolor{orange}{FAIL}\\
3. Prediction & Lower bound $>0$ in at least 3/4 cells; no 2-s failure; $<5\%$ explosive &
4/4; 4/4 positive at 2 s; 0\% explosive & \textcolor{green}{PASS}\\
4. Reproducibility & Median $S_{\rm sub}\geq0.80$ in at least 3/4; subset class agreement &
0/4; subset classes agree 4/4 but matching fails & \textcolor{orange}{FAIL}\\
5. Resolution & $R_{\rm track}<0.5$ in both recordings &
0.819 (ineligible) and 0.854 & \textcolor{orange}{FAIL}\\
\bottomrule
\end{longtable}

The frozen rule assigns \textbf{FAIL} if sparse/calcium-like known-system
recovery fails, even when empirical prediction passes. Gates 4 and 5 independently
reach the same practical conclusion. The coordinated-mode qualifier fails in
all four cells.

\begin{failbox}[Authorized conclusion]
The data contain forecastable temporal structure, but this pilot cannot identify
or track reproducible DMD modes at the proposed 300-frame resolution. Do not
claim awake/NREM/anaesthesia differences in DMD eigenvalues or subspaces from
this implementation, and do not proceed to intensive Grassmann tracking without
a newly reviewed measurement/window model.
\end{failbox}

\section{Limitations and defensible next options}

\subsection{What this pilot does not prove}

\begin{itemize}
  \item It does not prove that the underlying neural dynamics are non-linear or
  lack low-dimensional structure; it shows that the present observation/window
  pipeline cannot recover that structure reliably.
  \item It does not compare brain states statistically. Only two recordings and
  three fixed evaluation windows per cell were used.
  \item Bootstrap intervals condition on the selected windows and PCA basis.
  They quantify conditional operator-estimation perturbations, not uncertainty
  from test origins, window/neuron selection, preprocessing, model selection,
  animals, or experiments.
  \item The 1,500-frame diagnostic reselects ranks up to 32 and cannot isolate
  window length from model capacity.
  \item Stored smoothing can borrow support near label boundaries; the fixed
  edge trim reveals sensitivity but does not reconstruct the upstream smoother.
  \item Circular shifts use periodic wrap inside a legal bout. No shifted sample
  leaves the state/acquisition bout, but the wrap is a known surrogate artefact.
\end{itemize}

\subsection{If the project is revisited}

The most defensible new work would begin with a separate, prospectively frozen
plan and one of the following endpoints:

\begin{enumerate}
  \item a calcium-aware latent state-space/switching model fitted to fluorescence
  or deconvolution uncertainty, with simulation-calibrated identifiability;
  \item coarser event-count/rate observations and much longer windows, with the
  affine mean and diagonal baselines retained as primary comparators;
  \item static or slowly varying rate/covariance summaries if the scientific aim
  is brain-state characterization rather than modal dynamics;
  \item if modal tracking is retained, a prospectively eligible slow/predictive
  spectral cluster represented by ordered real-Schur subspaces, allowing
  real--complex transitions without post-hoc mode re-ranking; or
  \item additional recordings and animal-level validation before any state
  effect claim.
\end{enumerate}

The P arm may remain useful for a prediction benchmark, but the present results
do not justify using its eigenvectors as biological objects.

\section{Reproducibility and artifact inventory}

The final run contains 109 files (approximately 40 MB) across eight completed
stages. Generated results remain inside the private \texttt{DMD/} tree. The
accepted implementation-tree SHA-256 is:
\begin{center}
\small\texttt{95ee656eb4c9057fc004b76aebcd670418826be37bc620fc7afb470ebdca48f0}
\end{center}
It matches the current 24 hashed files under \texttt{DMD/src},
\texttt{DMD/scripts}, \texttt{DMD/tests}, and \texttt{DMD/configs}. The frozen
configuration hash is stored in \texttt{config\_frozen.json}; the Git revision
is:
\begin{center}
\small\texttt{ca0ddd9670b64bfbe8fd3565d005cb324642964d}
\end{center}
Because the DMD files are currently untracked, content hashes---not the Git
revision alone---identify the implementation.

Verification commands:
\begin{verbatim}
PYTHONPATH=DMD/src poetry run python -m unittest discover -s DMD/tests -v
poetry run ruff check DMD/src DMD/scripts DMD/tests
\end{verbatim}
All 25 tests pass and Ruff reports no violations. Tests cover selective HDF5
orientation, window legality, deterministic maximin placement, per-window B4,
frozen RMS scaling/PCA, known linear recovery, ridge scale invariance, exact
forecast powers, no cross-segment pairs, conjugate grouping, rejection of
full-space projectors, projector geometry, fixed-reference matching, no
bootstrap-block joins, within-bout circular shifts, stationary null sampling,
fixed split preservation, independent long-block diagonal selection, and
synchronized cell summaries.

Primary machine-readable entry points are:
\begin{itemize}
  \item \texttt{active\_run.json} --- accepted run pointer;
  \item \texttt{stage\_status.json} --- timestamps and stage summaries;
  \item \texttt{software\_manifest.json} --- versions and per-file hashes;
  \item \texttt{selection\_manifest.json} --- development-only winners;
  \item \texttt{predictive\_bootstrap\_cell\_summary.csv};
  \item \texttt{tracking\_stability\_cell\_summary.csv};
  \item \texttt{tracking\_resolution\_summary.csv};
  \item \texttt{negative\_control\_cell\_summary.csv}; and
  \item \texttt{final\_decision.json}.
\end{itemize}

\appendix
\section{Additional PC-trace diagnostics}

\begin{landscape}
\begin{figure}[p]
\centering
\includegraphics[width=0.98\linewidth,height=0.82\textheight,keepaspectratio]{\runroot/01_preprocessing/04_pc_traces__mouse02_sleep.png}
\caption{Mouse02 first four frozen-basis PC traces for every arm and state.}
\end{figure}
\end{landscape}

\begin{landscape}
\begin{figure}[p]
\centering
\includegraphics[width=0.98\linewidth,height=0.82\textheight,keepaspectratio]{\runroot/01_preprocessing/04_pc_traces__mouse07_ane.png}
\caption{Mouse07 first four frozen-basis PC traces for every arm and state.}
\end{figure}
\end{landscape}

\section{Exact selected window intervals}

All intervals below are zero-based and half-open:

\begin{longtable}{@{}llP{0.65\textwidth}@{}}
\toprule
Recording & Label & Chronological 300-frame windows\\
\midrule
\endfirsthead
\toprule
Recording & Label & Chronological 300-frame windows\\
\midrule
\endhead
mouse02 & awake &
D $[302,602)$; E $[2881,3181)$; D $[5962,6262)$; E $[9350,9650)$;
D $[11929,12229)$; E $[14508,14808)$\\
mouse02 & NREM &
D $[961,1261)$; E $[3897,4197)$; D $[7178,7478)$; E $[10270,10570)$;
D $[15464,15764)$; E $[18400,18700)$\\
mouse07 & awake &
D $[0,300)$; E $[1840,2140)$; D $[3680,3980)$; E $[5520,5820)$;
D $[7360,7660)$; E $[9200,9500)$\\
mouse07 & anaesthesia &
D $[9500,9800)$; E $[11340,11640)$; D $[13180,13480)$; E $[15020,15320)$;
D $[16860,17160)$; E $[18700,19000)$\\
\bottomrule
\end{longtable}

Long blocks are mouse02 awake $[12465,13965)$, mouse02 NREM
$[16918,18418)$, mouse07 awake $[4000,5500)$, and mouse07 anaesthesia
$[13500,15000)$.

\bibliographystyle{unsrtnat}
\bibliography{dmd_tracking_references}

\end{document}
